% =====================================================================
%  REVISÃO DA LITERATURA
%  Importe com \input{revisao-literatura} a partir do seu .tex principal.
%
%  Requer no preâmbulo:
%    \usepackage[utf8]{inputenc}   % se não estiver usando LuaLaTeX/XeLaTeX
%    \usepackage[brazil]{babel}
%
%  Bibliografia: as chaves \cite{} abaixo correspondem ao arquivo
%  referencias.bib fornecido junto com este capítulo.
%
%  Se a sua classe de documento for `article`, troque \chapter -> \section
%  e rebaixe os demais níveis.
% =====================================================================

\chapter{Revisão da Literatura}
\label{ch:revisao}

% ---------------------------------------------------------------------
\section{Redes de cápsulas}
\label{sec:revisao:capsulas}

Sabour, Frosst e Hinton propuseram as cápsulas em 2017~\cite{sabour2017dynamic}.
Cada cápsula emite um vetor: a norma indica a presença da entidade, a orientação
indica sua pose. As cápsulas de uma camada se acoplam às da seguinte por
roteamento por acordo, um processo iterativo, depois reformulado sobre cápsulas
matriciais~\cite{hinton2018matrix}.

Da pose explícita viriam duas vantagens. A primeira é robustez a pontos de vista
não vistos. A segunda é eficiência amostral: a rede não precisaria memorizar cada
transformação por augmentação ou por filtros redundantes.

Trabalhos posteriores baratearam a formulação. O DeepCaps aprofundou a
rede~\cite{rajasegaran2019deepcaps}; o Efficient-CapsNet chegou a 160 mil
parâmetros e trocou o roteamento iterativo por uma única passagem com
autoatenção~\cite{mazzia2021efficient}. É o Efficient-CapsNet que este trabalho
usa.

% ---------------------------------------------------------------------
\section{Uma controvérsia em aberto}
\label{sec:revisao:controversia}

\paragraph{Equivariância.}
O resultado favorável mais direto está em smallNORB~\cite{hinton2018matrix}.
Hinton e colegas igualaram cápsulas e CNN em pontos de vista familiares, depois
testaram as duas em azimutes e elevações excluídos do treino. O erro das cápsulas
ficou cerca de $30\%$ abaixo: $13{,}5\%$ contra $20{,}0\%$ em novos azimutes,
$12{,}3\%$ contra $17{,}8\%$ em novas elevações. O delineamento vale tanto quanto
o número --- igualar antes de medir isola o efeito do ponto de vista, e é a
lógica da razão de retenção usada aqui.

A evidência contrária é igualmente direta. Gu e Tresp não encontraram a
equivariância prometida sob transformações afins, e uma variante simplificada
subiu de $79\%$ para $93{,}21\%$ em AffNIST~\cite{gu2020affine}. Gu, Tresp e Hu
atribuíram a robustez das cápsulas à perda de margem e ao termo de reconstrução,
não aos componentes capsulares~\cite{gu2021notrobust}. Xi, Bing e Jin já haviam
medido queda acentuada em dados mais complexos que o MNIST~\cite{xi2017capsule}.
Mitterreiter e colegas vão além e apontam a árvore de análise dinâmica como a
razão de as cápsulas não escalarem~\cite{mitterreiter2023scale}.

\paragraph{Eficiência amostral.}
Jiménez-Sánchez, Albarqouni e Mateus compararam cápsulas e CNNs em imagem
médica~\cite{jimenez2018capsules}. As cápsulas igualaram ou superaram as CNNs com
menos dados anotados e toleraram melhor o desbalanceamento --- condição próxima à
do D-Fire. Essa reivindicação é a menos testada das duas: quase nunca é isolada
da augmentação, que sozinha cobre parte da variação que as cápsulas deveriam
dispensar.

% ---------------------------------------------------------------------
\section{Linhas de base e avaliação}
\label{sec:revisao:baselines}

As duas linhas de base vêm de paradigmas distintos: a ResNet-18,
convolucional~\cite{he2016resnet}, e o DeiT-Tiny, variante do \emph{Vision
Transformer}~\cite{dosovitskiy2021vit} treinável sem dados
externos~\cite{touvron2021deit}. A comparação entre esses dois paradigmas já deu
um alerta. Bai e colegas mostraram que a vantagem dos Transformers em robustez
desaparece quando protocolo, augmentação e parâmetros são
igualados~\cite{bai2021transformers}. Cápsulas tampouco são a única via para
equivariância --- redes equivariantes a grupos a obtêm por
construção~\cite{cohen2016group}.

A avaliação segue Hendrycks e Dietterich~\cite{hendrycks2019benchmarking}:
corrupções aplicadas ao teste, métricas normalizadas pela acurácia limpa. A
augmentação é o confundidor a controlar. RandAugment~\cite{cubuk2020randaugment}
e Random Erasing~\cite{zhong2020randomerasing} podem dar a redes comuns boa parte da
tolerância que se credita à arquitetura, então qualquer vantagem das cápsulas
precisa ser medida sob regimes de augmentação separados.

% ---------------------------------------------------------------------
\section{Lacuna}
\label{sec:revisao:lacuna}

Ficam três lacunas. A primeira é de arquitetura: os dois lados testaram a CapsNet
de 2017 ou cápsulas matriciais, não a formulação de passo único do
Efficient-CapsNet. A segunda é de domínio: os experimentos se concentram em
MNIST, AffNIST e smallNORB, e levantamentos recentes registram o
impasse~\cite{haq2023survey}. A terceira é de delineamento: robustez, augmentação
e fração de dados raramente aparecem cruzados, que é justamente o cruzamento
exigido pelo resultado de Bai e colegas~\cite{bai2021transformers}.

Este trabalho cobre as três. O delineamento cruza arquitetura, conjunto de dados
--- CIFAR-10~\cite{krizhevsky2009cifar} e D-Fire~\cite{venancio2022dfire} ---,
fração de treino, augmentação e semente, e avalia cada modelo sob transformações
ausentes do treino. O D-Fire reúne imagens de vigilância, de UAVs e da web, com
variação real de ponto de vista e escala. É a variação que interessa aqui:
revisões da área apontam a generalização a condições não vistas como o problema
central desses sistemas~\cite{elhanashi2025review}.

As hipóteses são falsificáveis e o desfecho está em aberto. Um resultado negativo
informa tanto quanto um positivo: nos dois casos a reivindicação passa por um
regime --- imagens naturais, protocolo igualado --- onde ainda não foi testada. O
protocolo está no Capítulo~\ref{ch:protocol}.
